% SPDX-FileCopyrightText: Copyright (c) 2024-2026 Yegor Bugayenko
% SPDX-License-Identifier: MIT

\section{Introduction}

Each object-oriented programming language offers its own set of objects,
  classes, types, functions, constants, traits, templates, and so on,
  which programmers can use off-the-shelf to model their specific use cases:
  ``Development Kit'' in Java~\citep{jdk2024,jdk8},
  ``Standard Library'' in
    C++~\citep{cpp2024,josuttis2012cpp},
    Python~\citep{python2024,hellmann2017python},
    Swift~\citep{swift2024,deitel2015swift},
    Rust~\citep{rust2024,blandy2021rust},
    and
    Ruby,
  ``.NET Standard'' in C\#~\citep{net2023,abrams2005net},
  and ``Built-in Objects'' in JavaScript~\citep{js2024,crockford2008js}.

Each standard library (SL) freely defines its own principles of abstraction.
For example, to get the absolute value of a numeric object \ff{x} in Java,
  one has to call a static method of another class \ff{Math.abs(x)},
  while in Ruby it would be a property of the same object \ff{x.abs}.
For example, in Java, an attempt to get a value from a hash map by
  a key will produce \ff{null} in case of its absence,
  while in C++ it will return an empty iterator.
There are many similar examples demonstrating
  the absence of a common ground between SLs.

Earlier, \citet{booch1990design} suggested their own components for OOP.
We suggest our own taxonomy of objects for \eolang{}\footnote{%
  \url{https://www.eolang.org}, 9.9.9},
  a strictly formal~\citep{kudasov2022formalizing} object-oriented
  programming language with an intentionally reduced
  feature set~\citep{bugayenko2021eolang}.

An object in \eolang{} is either a composition of other objects
  or an atom, which is a foreign function interface
  to a lower-level runtime such as, for example, JVM or Assembly.
Throughout the paper, we use dark blue color
  to highlight the names of objects,
  as opposed to other pieces of fixed-width text.
We use orange color to highlight the names of atoms.

In this paper%
  \footnote{\raggedright%
    \LaTeX{} sources of this paper are maintained in the
    \href{https://github.com/REPOSITORY}{REPOSITORY} GitHub repository,
    the rendered version is
    \href{https://github.com/REPOSITORY/releases/tag/0.0.0}{\ff{0.0.0}}.}
  we don't provide exact specification of each object,
  in order to avoid conflicts with their exact definitions
  in the Objectionary\footnote{\url{https://github.com/objectionary/home}}.
